\section{Resource Mapping}\label{sec:methodology}

We compare forty-one selected evaluation resources to see how they use the
three components of the truth contract: the task oracle, permission scope,
and required status marking. We also examine how they evaluate style.
In particular, we ask whether they compare different wordings of a response
while keeping claim content and observed status marking fixed. This
comparison, which we call the \emph{resource mapping}, describes how the
resources evaluate responses. It does not test whether our labeling rule
works in practice.

\subsection{Selection and Scope}

An \emph{evaluation resource} is a benchmark, dataset, or protocol that
specifies how responses are evaluated. We selected resources for different
tasks, from factuality evaluation to creative writing and problem solving.
The aim is to compare how these tasks are evaluated. The selection does not
cover all existing resources or show how common each approach is across
the field. Tables~\ref{tab:full-mapping-factual-a}--\ref{tab:full-mapping-creative}
list the 41 resources and their references.

The available records do not explain how the resources were first found
and selected. This information is needed for others to repeat the selection.
We also could not examine every source in the same detail: for some resources,
we read the methods; for others, we relied on the authors' summaries; and
one source was only partly accessible. If we did not identify a criterion
in these materials, this does not prove that the resource lacks it.
The recent resources discussed in Section~\ref{sec:related} remain outside
this comparison until they have been examined in the same way.
% AUTHOR ACTION: document the actual resource search and selection procedure.
% Do not invent dates, queries, or inclusion criteria.

\subsection{Comparison Criteria}

For each resource, we examine three groups of information:
\begin{itemize}
\item \textbf{Truth contract.} What reference evidence or checks does the task
oracle use? What content does the task permit? How must that content be marked?
\item \textbf{Task.} What does the response have to do, and how are its quality
or success evaluated?
\item \textbf{Style.} Does the evaluation check the requested style? Does it
compare different wordings with the same claim content and observed status
marking?
\end{itemize}
These questions use the distinctions introduced in
Section~\ref{sec:framework}. A check that a solution works does not necessarily
establish whether every claim in its explanation is supported.

We group the resources into six task categories. The names and resource IDs
are the same in the tables and Figure~\ref{fig:profiles}:
\begin{itemize}
\item \textbf{Factuality evaluation} (IDs 1--22) checks answers or summaries
against reference evidence.
\item \textbf{Intent evaluation} (23) checks whether a response correctly
interprets the user's request.
\item \textbf{Mixed-task and tool-use evaluation} (24--27) combines different
tasks or requires the use of tools.
\item \textbf{Creative writing and reasoning} (28--32, 35--36) evaluates
stories, written explanations, or different ways of solving a problem.
\item \textbf{Divergent thinking and linguistic creativity} (33--34)
evaluates the variety and originality of ideas or language.
\item \textbf{Constrained problem solving} (37--41) checks whether a solution
meets the task's requirements, such as working with limited resources or
passing code tests.
\end{itemize}
These categories describe the tasks. They do not rank the resources or
determine whether an individual claim is permitted. They reflect our reading
of the resources and have not been checked through independent annotation.

\subsection{Resource Tables}\label{subsec:resource-tables}

Each row describes one resource. The columns give its name and reference,
the task and the evidence or checks used, other evaluation criteria, and
our reading of what content the task allows.

``No extra quality measure identified'' means that we did not identify a
quality measure beyond the evidence or correctness check already listed.
``Only supported content (inferred)'' means that we read the task as requiring
claims to be supported. For creative tasks, the final column describes what
we understand the instructions to allow. When it says ``Permission not
established by this review,'' we have not established which unsupported
additions are permitted. This does not mean that all additions are forbidden.
Required status marking and style are discussed after the tables.

\begingroup
\small
\setlength{\tabcolsep}{3pt}
\renewcommand{\arraystretch}{1.0}
\begin{longtable}{@{}P{.20\linewidth}P{.32\linewidth}P{.24\linewidth}P{.19\linewidth}@{}}
\caption{\textbf{Factuality evaluation (resources 1--13).} The final column gives our reading of the task instructions. Permission for individual claims still needs to be checked.}
\label{tab:full-mapping-factual-a}\\
\toprule
\textbf{Resource} & \textbf{What is checked, and against what?} & \textbf{Other evaluation criteria} & \textbf{Allowed content: our reading}\\\midrule
\endfirsthead
\multicolumn{4}{l}{\textbf{Table \thetable{} (continued)}}\\
\toprule
\textbf{Resource} & \textbf{What is checked, and against what?} & \textbf{Other evaluation criteria} & \textbf{Allowed content: our reading}\\\midrule
\endhead
\midrule\multicolumn{4}{r}{\emph{Continued on next page}}\\
\endfoot
\bottomrule
\endlastfoot
\multicolumn{4}{@{}p{.96\linewidth}@{}}{\textbf{Task family: factuality evaluation}}\\*
1.~FEVER \citep{thorne_fever_2018} & Fact checking\newline Wikipedia evidence sentences. & No extra quality measure identified & Only supported content (inferred)\\\addlinespace[3pt]
2.~TruthfulQA \citep{lin_truthfulqa_2022} & Factual QA\newline Sourced reference answers; human adjudication for generations. & Informativeness & Only supported content (inferred)\\\addlinespace[3pt]
3.~PopQA \citep{mallen_when_2023} & Long-tail QA\newline Wikidata. & No extra quality measure identified & Only supported content (inferred)\\\addlinespace[3pt]
4.~SimpleQA \citep{wei_measuring_2024} & Factual QA\newline Verified reference answers. & Answer correctness and not-attempted outcome & Only supported content (inferred)\\\addlinespace[3pt]
5.~TriviaQA \citep{joshi_triviaqa_2017} & Reading QA\newline Gold answers plus Web/Wikipedia documents. & No extra quality measure identified & Only supported content (inferred)\\\addlinespace[3pt]
6.~DefAn \citep{rahman_defan_2024} & Factual QA\newline Definitive reference answers. & Prompt alignment and cross-run consistency & Only supported content (inferred)\\\addlinespace[3pt]
7.~HaluEval \citep{li_halueval_2023} & Multi-task hallucination recognition\newline Task references or source text; Web-assisted adjudication. & Specificity and relevance folded into some labels & Only supported content (inferred)\\\addlinespace[3pt]
8.~FactBench \citep{bayat_factbench_2024} & Open-ended factuality evaluation\newline Web-retrieved material. & No extra quality measure identified & Only supported content (inferred)\\\addlinespace[3pt]
9.~HalluLens \citep{bang_hallulens_2025} & Extrinsic/intrinsic hallucination benchmark\newline Wikipedia, nonexistent entities, or task context. & Correctness and refusal/acceptance & Only supported content (inferred)\\\addlinespace[3pt]
10.~FRANK \citep{pagnoni_understanding_2021} & Summarization\newline Source documents. & No extra quality measure identified & Only supported content (inferred)\\\addlinespace[3pt]
11.~QAGS \citep{wang_asking_2020} & Summarization\newline Source documents. & No extra quality measure identified & Only supported content (inferred)\\\addlinespace[3pt]
12.~SummEval \citep{fabbri_summeval_2021} & Summarization\newline Source documents. & Automatic metrics and human summarization-quality judgments & Only supported content (inferred)\\\addlinespace[3pt]
13.~XSum \citep{narayan_xsum_2018} & Summarization\newline BBC source article. & ROUGE, human preference, and information preservation & Only supported content (inferred)\\
\end{longtable}
\endgroup

\begingroup
\small
\setlength{\tabcolsep}{3pt}
\renewcommand{\arraystretch}{1.0}
\begin{longtable}{@{}P{.20\linewidth}P{.32\linewidth}P{.24\linewidth}P{.19\linewidth}@{}}
\caption{\textbf{Factuality evaluation, intent evaluation, and mixed-task and tool-use evaluation (resources 14--27).} The final column gives our reading of the task instructions. Permission for individual claims still needs to be checked.}
\label{tab:full-mapping-factual-b}\\
\toprule
\textbf{Resource} & \textbf{What is checked, and against what?} & \textbf{Other evaluation criteria} & \textbf{Allowed content: our reading}\\\midrule
\endfirsthead
\multicolumn{4}{l}{\textbf{Table \thetable{} (continued)}}\\
\toprule
\textbf{Resource} & \textbf{What is checked, and against what?} & \textbf{Other evaluation criteria} & \textbf{Allowed content: our reading}\\\midrule
\endhead
\midrule\multicolumn{4}{r}{\emph{Continued on next page}}\\
\endfoot
\bottomrule
\endlastfoot
\multicolumn{4}{@{}p{.96\linewidth}@{}}{\textbf{Task family: factuality evaluation (continued)}}\\*
14.~CNN/\allowbreak DailyMail \citep{hermann_teaching_2015,wang_asking_2020} & News corpus used in summarization evaluation\newline News source article. & No extra quality measure identified & Only supported content (inferred)\\\addlinespace[3pt]
15.~DelucionQA \citep{sadat_delucionqa_2023} & Domain QA\newline Retrieved chunks from one Jeep Gladiator manual. & Answerability, coverage, and abstention & Only supported content (inferred)\\\addlinespace[3pt]
16.~HalluQA \citep{cheng_evaluating_2023} & Chinese adversarial QA\newline Answer sets backed by sources or explanations. & Answer correctness and refusal handling & Only supported content (inferred)\\\addlinespace[3pt]
17.~CR-LLM \citep{li_cr-llm_2024} & Concept reasoning\newline Reference concepts plus GPT-4 contextual adjudication. & Concept granularity & Only supported content (inferred)\\\addlinespace[3pt]
18.~HotpotQA \citep{yang_hotpotqa_2018} & Multi-hop QA\newline Wikipedia evidence passages. & No extra quality measure identified & Only supported content (inferred)\\\addlinespace[3pt]
19.~HalluRAG \citep{ridder_hallurag_2025} & RAG hallucination detection\newline Provided passages/answers and derived sentence labels. & No extra quality measure identified & Only supported content (inferred)\\\addlinespace[3pt]
20.~RealTime QA \citep{kasai_realtime_2024} & Real-time QA\newline Time-indexed reference answers. & No extra quality measure identified & Only supported content (inferred)\\\addlinespace[3pt]
21.~HypoTermQA \citep{uluoglakci_hypotermqa_2024} & Hypothetical-term evaluation\newline Term-validity labels and Wikipedia definitions. & No extra quality measure identified & Only supported content (inferred)\\\addlinespace[3pt]
22.~AA-Omniscience \citep{jackson_aaomniscience_2025,artificialanalysis_aaomniscience_public_2025} & Cross-domain factual QA\newline Source-derived reference answers, evaluated by an LLM grader. & Accuracy, incorrect-answer penalty, and abstention-sensitive reliability & Only supported content (inferred)\\\addlinespace[3pt]
\midrule
\multicolumn{4}{@{}p{.96\linewidth}@{}}{\textbf{Task family: intent evaluation}}\\*
23.~FaithQA \citep{hao_faithqa_2025} & Intent verification\newline RAG reports; Wikipedia for factual QA; none for creative writing. & Query-constraint satisfaction & Permission not established by this review\\\addlinespace[3pt]
\midrule
\multicolumn{4}{@{}p{.96\linewidth}@{}}{\textbf{Task family: mixed-task and tool-use evaluation}}\\*
24.~BIG-Bench \citep{srivastava_beyond_2023} & Multi-task reasoning\newline Subtask targets and adjudication. & No extra quality measure identified & Permission not established by this review\\\addlinespace[3pt]
25.~BIG-Bench Hard \citep{suzgun_challenging_2023} & Reasoning suite\newline Task-specific reference answers. & No extra quality measure identified & Permission not established by this review\\\addlinespace[3pt]
26.~ToolQA \citep{zhuang_toolqa_2023} & Tool-use QA\newline Fixed corpora and generated gold answers. & No extra quality measure identified & Permission not established by this review\\\addlinespace[3pt]
27.~HIC-Bench \citep{yang_hicbench_2025} & Hallucination typing\newline CDKB/Wikipedia principles and human-verified judgments. & Originality, feasibility, value, and composite metrics & Permission not established by this review\\
\end{longtable}
\endgroup

\begingroup
\small
\setlength{\tabcolsep}{3pt}
\renewcommand{\arraystretch}{1.0}
\begin{longtable}{@{}P{.20\linewidth}P{.32\linewidth}P{.24\linewidth}P{.19\linewidth}@{}}
\caption{\textbf{Creative writing and reasoning, divergent thinking and linguistic creativity, and constrained problem solving (resources 28--41).} The final column gives our reading of the task instructions. Permission for individual claims still needs to be checked.}
\label{tab:full-mapping-creative}\\
\toprule
\textbf{Resource} & \textbf{What is checked, and against what?} & \textbf{Other evaluation criteria} & \textbf{Allowed content: our reading}\\\midrule
\endfirsthead
\multicolumn{4}{l}{\textbf{Table \thetable{} (continued)}}\\
\toprule
\textbf{Resource} & \textbf{What is checked, and against what?} & \textbf{Other evaluation criteria} & \textbf{Allowed content: our reading}\\\midrule
\endhead
\midrule\multicolumn{4}{r}{\emph{Continued on next page}}\\
\endfoot
\bottomrule
\endlastfoot
\multicolumn{4}{@{}p{.96\linewidth}@{}}{\textbf{Task family: creative writing and reasoning}}\\*
28.~CreativityPrism \citep{hou_creativityprism_2025} & Multi-task creativity\newline Tests or correctness judges where applicable. & Task-specific rubrics & Invention within task instructions\\\addlinespace[3pt]
29.~LitBench \citep{fein_litbench_2025} & Story evaluation\newline No external factual reference identified (human preferences assess quality). & Human preferences & Fictional content within the story\\\addlinespace[3pt]
30.~Short Story \citep{ismayilzada_evaluating_2025} & Story generation\newline No external factual reference identified (cue words constrain the task, not claim truth). & Novelty, diversity, complexity, and judged creativity & Fictional content within the story\\\addlinespace[3pt]
31.~WritingBench \citep{wu_writingbench_2025} & Generative writing\newline Provided materials when present. & Writing quality and requested style & Invention within task instructions\\\addlinespace[3pt]
32.~Art-or-Artifice / TTCW \citep{chakrabarty_art_2024} & Creative writing\newline No external factual reference identified (expert tests assess creativity). & Expert rubric & Fictional content within the story\\\addlinespace[3pt]
35.~Automated Creativity \citep{li_automated_2025} & Creative writing\newline No external factual reference identified (reference texts anchor quality, not claim truth). & Reference-text quality anchor & Invention within task instructions\\\addlinespace[3pt]
36.~Hallucinating LLM Could Be Creative \citep{sun_hallucinating_2024} & Reasoning-path evaluation\newline Task correctness; consistency/clustering are diagnostics. & Correctness, consistency, diversity & Permission only partly characterized\\\addlinespace[3pt]
\midrule
\multicolumn{4}{@{}p{.96\linewidth}@{}}{\textbf{Task family: divergent thinking and linguistic creativity}}\\*
33.~Automated DT scoring \citep{organisciak_beyond_2023} & Divergent thinking\newline No external factual reference identified. & Novelty and variety of ideas & New ideas within the task\\\addlinespace[3pt]
34.~Verbal Creativity \citep{dinu_testing_2025} & Linguistic creativity\newline No external factual reference identified. & Originality, flexibility, elaboration & Creative language within the task\\\addlinespace[3pt]
\midrule
\multicolumn{4}{@{}p{.96\linewidth}@{}}{\textbf{Task family: constrained problem solving}}\\*
37.~MacGyver \citep{tian_macgyver_2025} & Problem solving\newline Human judgment of feasibility and efficiency. & Feasibility and efficiency & Permission not established by this review\\\addlinespace[3pt]
38.~EscapeBench \citep{qian_escapebench_2025} & Agent puzzles\newline Game state and progress. & Feasibility or correctness & Permission not established by this review\\\addlinespace[3pt]
39.~CreativEval \citep{delorenzo_creativeval_2024} & Code generation\newline Functionality under Verilog testbenches. & Fluency, flexibility, originality, elaboration & Permission not established by this review\\\addlinespace[3pt]
40.~NeoCoder \citep{lu_benchmarking_2025} & Code generation\newline Test success and constraint adherence. & Correctness, constraint adherence, divergence & Permission not established by this review\\\addlinespace[3pt]
41.~CreativeMath \citep{ye_assessing_2025} & Mathematical solution generation\newline Reference solutions and unanimous LLM correctness judgment. & Correctness and solution novelty & Permission not established by this review\\
\end{longtable}
\endgroup

\subsection{Resource Analysis Criteria}\label{sec:mapping}

\paragraph{Task oracle}
The resources use different reference evidence and checks. FEVER uses
evidence sentences (ID~1), QAGS compares answers obtained from a summary
and its source (11), and SimpleQA checks answers against reference answers
(4) \citep{thorne_fever_2018,wang_asking_2020,wei_measuring_2024}.
The reference matters because the same claim can be supported by one
reference and left unknown by another, as Example~A in
Section~\ref{sec:examples} illustrates.

For constrained problem solving, the checks ask whether a solution works.
MacGyver uses human judgments of feasibility and efficiency (37);
EscapeBench checks progress in a game (38); CreativEval uses Verilog tests
(39); NeoCoder checks test success and task constraints (40); and
CreativeMath uses LLM judgments of mathematical correctness (41)
\citep{tian_macgyver_2025,qian_escapebench_2025,delorenzo_creativeval_2024,
lu_benchmarking_2025,ye_assessing_2025}. These checks can evaluate a solution
without verifying every claim in its explanation. The tables describe
what each check does; they do not establish that it performs every part
of the task oracle defined in Section~\ref{subsec:contract}.

For six resources, we did not identify an external factual reference in the
evaluation criteria we examined: LitBench, Short Story, Art-or-Artifice,
Automated DT scoring, Verbal Creativity, and Automated Creativity
(IDs 29, 30, 32--35). Their evaluations still use task instructions,
requirements for stories, reference texts for quality, or creativity
criteria. Our reading does not establish that these resources never check
factual claims.

\paragraph{Permission scope}
For the 22 resources in factuality evaluation (IDs 1--22), we read the tasks
as allowing supported content only. For eight resources (28--35), we infer
what content is allowed from the task instructions or the fictional setting.
For resource 36, we have identified only part of the permission conditions.
For the remaining ten resources (23--27 and 37--41), we have not established
which unsupported additions are allowed.

These descriptions explain our reading of the tasks. We still need to check
whether their instructions cover the full permission scope defined in
Section~\ref{subsec:contract}. Section~\ref{sec:agenda} describes this work.

\paragraph{Required status marking}
For these resources, we have not yet established both what status marking
the task requires and how the evaluation checks it. We therefore do not
classify resources as having or lacking a marking requirement. A requirement
may be present in instructions that our review has not fully examined.

Some resources evaluate whether a response answers or declines to answer.
SimpleQA includes a not-attempted outcome, HalluLens evaluates refusal, and
DelucionQA records answerability and whether a response answers
\citep{wei_measuring_2024,bang_hallulens_2025,sadat_delucionqa_2023}.
AA-Omniscience also includes abstention in its reliability evaluation
\citep{jackson_aaomniscience_2025}. This differs from checking how a claim
is marked: declining to answer may produce no claim about the answer at all,
whereas a permitted hypothesis introduces a claim whose status must be clear.

\subsection{Joint Evaluation and Style Variation}

Style evaluation can ask two questions. The first is whether a response
uses the requested style. WritingBench includes this in its evaluation
\citep{wu_writingbench_2025}. The second is whether changing the wording
changes an evaluator's judgment when claim content and observed status
marking remain fixed. We have not documented this second comparison in the
selected resources. This does not establish that it is absent from every
task they include or from other resources. For actions or executable code,
a comparison of prose style may not apply.

Figure~\ref{fig:profiles} groups the resources by task category. Its center
shows what we want to evaluate together: creativity and divergent thinking,
and whether claims are licensed divergence or hallucination under a truth
contract. Some resources already combine creativity with correctness checks
or hallucination evaluation, including HIC-Bench and resources for
constrained problem solving in Table~\ref{tab:full-mapping-creative}.
We still need to examine whether they use reference evidence, permission
scope, and required status marking to distinguish permitted content from
hallucination. The figure states this question; it does not prove that no
existing resource addresses it.

Style remains a separate question. An evaluation can check creative quality
and factuality without testing whether a change in wording affects its
judgment. The next section uses the same contract components to explain
what should change, or stay the same, in five constructed examples.

\begin{figure}[!ht]
\centering
\resizebox{\textwidth}{!}{\input{figures/resource_clusters.tex}}
\caption{\textbf{Task categories and the evaluation proposed in this paper.}
The boxes use the same six categories as the resource tables in
Section~\ref{subsec:resource-tables}. Names and IDs identify the resources
in each group. The center brings together creativity, divergent thinking,
and hallucination evaluation under a truth contract. Some resources already
combine these questions; the figure does not establish that this combination
is absent from prior work. Box positions and sizes do not give scores or
rank resources. Tables~\ref{tab:full-mapping-factual-a}--\ref{tab:full-mapping-creative}
give the details for each resource.}
\label{fig:profiles}
\end{figure}
\FloatBarrier

